\documentclass{article}

\usepackage[margin=1in]{geometry}
\usepackage{amsmath,amsthm,amssymb}
\usepackage[spanish]{babel}
\usepackage{enumitem}

\theoremstyle{definition}
\newtheorem{definition}{Definición}[section]

% Number sets
\newcommand{\R}{\mathbb{R}}
\newcommand{\Z}{\mathbb{Z}}
\newcommand{\N}{\mathbb{N}}
\newcommand{\Q}{\mathbb{Q}}
\newcommand{\C}{\mathbb{C}}

% Topology operators
\newcommand{\cl}[1]{\overline{#1}}              % Closure
\newcommand{\Int}[1]{{#1}^\circ}                % Interior
\newcommand{\bd}{\partial}                       % Boundary
\newcommand{\Ent}[1]{\mathcal{O}(#1)}           % Neighborhoods (Entornos)
\newcommand{\pf}{\mathcal{P}_F}                 % Finite parts
\newcommand{\partes}{\mathcal{P}}               % Power set

% Useful shortcuts
\newcommand{\sii}{\iff}                          % Si y sólo si
\newcommand{\imp}{\implies}                      % Implies
\newcommand{\setdiff}{\smallsetminus}            % Set difference
\newcommand{\coint}[2]{\left[ #1, #2 \right)}

% Net convergence notation
\newcommand{\evto}[1]{\stackrel{E}{\hookrightarrow} #1}
\newcommand{\frto}[1]{\stackrel{F}{\hookrightarrow} #1}

% Exercise environment
\newenvironment{ejercicio}[1]{\begin{trivlist}
\item[\hskip \labelsep {\bfseries Ejercicio}\hskip \labelsep {\bfseries #1.}]}{\end{trivlist}}

\setlist[enumerate,1]{label=(\alph*), leftmargin=*, itemsep=2pt}

\begin{document}

\title{Topología: Práctica de Repaso\\Primer Parcial}
\author{Bustos Jordi\\Espacios Topológicos, Conexidad y Redes}

\maketitle

\section{Espacios Topológicos y Separación}

\begin{ejercicio}{1}
    Sea \((X, \leq)\) un conjunto totalmente ordenado. Se define la \textit{topología del orden} \(\tau_\leq \) como la generada por los intervalos abiertos \(\{ (a,b) : a, b \in X,\, a < b \}\), donde \((a,b) = \{ x \in X : a < x < b \}\).
    Notar que si \(X\) tiene un mínimo \(m\) o un máximo \(M\), entonces también se incluyen los intervalos de la forma \(\coint{m}{b}\) o \((a, M]\) respectivamente. Demostrar las siguientes afirmaciones sobre \((X, \tau_\leq)\):
    \begin{enumerate}
        \item Demostrar que \((X, \tau_\leq)\) es \(T_2\).
        \item Mostrar que si \(X = \R \) con el orden usual, entonces \(\tau_\leq \) coincide con la topología usual de \(\R \).
        \item Describir la clausura del intervalo abierto \((a,b)\) en la topología del orden sobre \(\R \), y dar un ejemplo de un conjunto totalmente ordenado en el que \(\cl{(a,b)} \neq [a,b]\).
    \end{enumerate}
\end{ejercicio}

\begin{ejercicio}{2}
    Sean \(\tau_1 \subseteq \tau_2\) dos topologías sobre un mismo conjunto \(X\). Analizar las siguientes implicaciones; demostrarlas o dar un contraejemplo:
    \begin{enumerate}
        \item Si \((X, \tau_2)\) es \(T_1\), entonces \((X, \tau_1)\) es \(T_1\).
        \item Si \((X, \tau_1)\) es \(T_2\), entonces \((X, \tau_2)\) es \(T_2\).
        \item Si \((X, \tau_1)\) es normal, entonces \((X, \tau_2)\) es normal.
    \end{enumerate}
\end{ejercicio}

\begin{ejercicio}{3}
    Sea \((X, \tau)\) un ET y \(A \subseteq X\) un subespacio. Demostrar que:
    \begin{enumerate}
        \item Si \((X, \tau)\) es \(T_1\), entonces \((A, \tau_A)\) es \(T_1\).
        \item Si \((X, \tau)\) es \(N_1\), entonces \((A, \tau_A)\) es \(N_1\).
        \item Si \((X, \tau)\) es \(N_2\), entonces \((A, \tau_A)\) es \(N_2\).
    \end{enumerate}
\end{ejercicio}

\begin{ejercicio}{4}
    Sea \((X, \tau)\) un ET que satisface el primer axioma de numerabilidad \((N_1)\).
    \begin{enumerate}
        \item Demostrar que para todo \(A \subseteq X\) y todo \(x \in X\): \( x \in \cl{A} \) si y sólo si existe una sucesión \({(a_n)}_{n \in \N}\) en \(A\) tal que \(a_n \to x\). (Nota: en espacios que no son \(N_1\) esto puede fallar.)
        \item Concluir que en un ET \(N_1\), un subconjunto \(A\) es cerrado si y sólo si es \textit{secuencialmente cerrado}: toda sucesión en \(A\) convergente en \(X\) tiene límite en \(A\).
    \end{enumerate}
\end{ejercicio}

\clearpage

\section{Conexidad}

\begin{ejercicio}{5}
    Sea \(A\) un subespacio conexo de \((X, \tau)\). Demostrar que \(\cl{A}\) es conexo. Más aún, demostrar que si \(A \subseteq B \subseteq \cl{A}\), entonces \(B\) también es conexo.
\end{ejercicio}

\begin{ejercicio}{6}
    Sea \(f : X \to Y\) continua.
    \begin{enumerate}
        \item Si \(X\) es conexo, demostrar que \(f(X)\) es conexo.
        \item Si \(f\) es además abierta y \(X\) es localmente conexo, demostrar que \(f(X)\) es localmente conexo.
        \item Deducir que la conexidad y la conexidad local son invariantes topológicos (se preservan bajo homeomorfismos).
    \end{enumerate}
\end{ejercicio}

\begin{ejercicio}{7}
    Sea \(f : X \to \R \) continua con \(X\) conexo.
    \begin{enumerate}
        \item Demostrar que \(f(X)\) es un intervalo (posiblemente no acotado o degenerado).
        \item Deducir el \textbf{Teorema del Valor Intermedio}: si \(f : [a, b] \to \R \) es continua, \(f(a) < c < f(b)\), entonces existe \(x_0 \in (a, b)\) tal que \(f(x_0) = c\).
        \item Probar que todo polinomio de grado impar \(p : \R \to \R \) tiene al menos una raíz real.
    \end{enumerate}
\end{ejercicio}

\begin{ejercicio}{8}
    Sea \((X, \tau)\) un ET conexo y \(A \subsetneq X\) un subconjunto propio y no vacío.
    \begin{enumerate}
        \item Demostrar que \(\bd A \neq \varnothing \).
        \item Deducir que los únicos subconjuntos de \(X\) que son simultáneamente abiertos y cerrados son \(\varnothing \) y \(X\).
        \item Sea \(f : X \to \{0, 1\} \) continua (donde \( \{0,1\} \) lleva la topología discreta). Demostrar que \(f\) es constante.
    \end{enumerate}
\end{ejercicio}

\clearpage

\section{Redes y Convergencia}

\begin{ejercicio}{9}
    Sea \((X, \tau)\) un ET, \(A \subseteq X\) y \(x \in X\). Demostrar que
    \[
        x \in \cl{A} \iff \exists \text{ red } {(x_i)}_{i \in I} \text{ en } A \text{ tal que } x_i \to x.
    \]
    Deducir que \(A\) es cerrado si y sólo si es cerrado bajo límites de redes.
\end{ejercicio}

\begin{ejercicio}{10}
    Sea \((X, \tau)\) un ET Hausdorff. Demostrar que toda red en \(X\) tiene a lo sumo un límite. Analizar si la recíproca es cierta: si toda red tiene a lo sumo un límite, ¿es \(X\) necesariamente \(T_2\)?
\end{ejercicio}

\begin{ejercicio}{11}
    Sea \(X\) un conjunto infinito dotado de la topología cofinita \(\tau_{CF}\).
    \begin{enumerate}
        \item Demostrar que si \({(x_i)}_{i \in I}\) es una red en \(X\) que para todo subconjunto finito \(F \subsetneq X\) está eventualmente en \(X \setminus F\), entonces \(x_i \to p\) para todo \(p \in X\).
        \item Sea \({(x_n)}_{n \in \N}\) una sucesión en \(X\) con todos sus términos distintos. Demostrar que converge a todo punto de \(X\) en la topología \(\tau_{CF}\). Deducir que en \((\R, \tau_{CF})\) la convergencia de sucesiones no caracteriza la topología (compare con el ejercicio anterior de \(N_1\)).
        \item Demostrar que toda función continua \(f : (\R, \tau_{CF}) \to (\R, \tau_{usual})\) es constante.
    \end{enumerate}
\end{ejercicio}

\begin{ejercicio}{12}
    Sea \((X, \tau)\) un ET, \({(x_i)}_{i \in I}\) una red en \(X\) y \(x \in X\). Decimos que \(x\) es \textit{punto de acumulación} de la red si para todo entorno \(U \in \Ent{x}\) y todo \(i_0 \in I\) existe \(i \geq i_0\) tal que \(x_i \in U\).
    \begin{enumerate}
        \item Demostrar que si \(x_i \to x\), entonces \(x\) es punto de acumulación de la red.
        \item Demostrar que \(x\) es punto de acumulación de \({(x_i)}_{i \in I}\) si y sólo si existe una subred de \({(x_i)}\) que converge a \(x\).
        \item Deducir que en un espacio \(T_2\), si una red tiene un único punto de acumulación \(x\) y converge a algún punto, entonces converge a \(x\).
    \end{enumerate}
\end{ejercicio}

\clearpage

\subsection*{Ejercicios Complementarios}
\begin{ejercicio}{C1}
    Sea \((X, \tau)\) un ET y \(U \subseteq X\). Probar que \(U\) es abierto si y sólo si para toda red \({(x_i)}_{i \in I}\) en \(X\) tal que \(x_i \to x \in U\), la red está eventualmente en \(U\).
\end{ejercicio}

\begin{ejercicio}{C2}
    Sea \(Y\) un subespacio cerrado de un ET normal \((X, \tau)\). Demostrar que \((Y, \tau_Y)\) es normal.
\end{ejercicio}

\begin{ejercicio}{C3}
    Sea \((X, d)\) un EM completo y sea \({(F_n)}_{n \geq 1}\) una sucesión decreciente de subconjuntos cerrados no vacíos de \(X\) tal que \(\mathrm{diam}(F_n) \to 0\). Demostrar que \(\bigcap_{n \geq 1} F_n\) es no vacío y tiene exactamente un punto. (Teorema de Cantor de intersección anidada.)
\end{ejercicio}

\clearpage

\section{Segunda Mitad: Continuidad, Compactos y Compacidad Local}

\begin{ejercicio}{13}
    Sean \(X\) e \(Y\) espacios topológicos con \(Y\) Hausdorff, \(D \subseteq X\) denso y \(f, g : X \to Y\) continuas.
    \begin{enumerate}
        \item Demostrar que el conjunto \(A = \{x \in X : f(x) = g(x)\}\) es cerrado en \(X\).
        \item Si \(f|_D = g|_D\), deducir que \(f = g\) en todo \(X\).
        \item Dar un contraejemplo que muestre que la conclusión puede fallar si \(Y\) no es Hausdorff.
    \end{enumerate}
\end{ejercicio}

\begin{ejercicio}{14}
    Sea \(X\) un espacio topológico y para cada \(n \in \N\) sea \(f_n : X \to \R\) una función continua. Definimos
    \[
        F : X \to \R^{\mathbb{N}}, \qquad F(x) = (f_1(x), f_2(x), \dots).
    \]
    \begin{enumerate}
        \item Demostrar que \(F\) es continua cuando \(\R^{\mathbb{N}}\) lleva la topología producto.
        \item Aplicar el punto anterior al caso \(f_n = f\) para una misma función continua \(f : X \to \R\).
        \item Analizar si la conclusión del punto (a) sigue siendo cierta si en \(\R^{\mathbb{N}}\) se considera la topología caja. En caso negativo, dar un contraejemplo explícito.
    \end{enumerate}
\end{ejercicio}

\begin{ejercicio}{15}
    Sea \(X\) un espacio topológico Hausdorff y \(K \subseteq X\) compacto.
    \begin{enumerate}
        \item Demostrar que para todo \(x \in X \setminus K\) existen abiertos disjuntos \(U, V \subseteq X\) tales que \(x \in U\) y \(K \subseteq V\).
        \item Deducir que todo compacto de \(X\) es cerrado.
        \item Sea \(f : X \to Y\) continua y biyectiva, donde \(X\) es compacto e \(Y\) es Hausdorff. Probar que \(f\) es un homeomorfismo.
    \end{enumerate}
\end{ejercicio}

\begin{ejercicio}{16}
    Sea \(X\) un espacio topológico compacto y \(\mathcal{F}\) una familia de cerrados de \(X\).
    \begin{enumerate}
        \item Demostrar que si \(\mathcal{F}\) tiene la propiedad de intersección finita, entonces \(\bigcap_{F \in \mathcal{F}} F \neq \varnothing\).
        \item Deducir que toda sucesión decreciente \(K_1 \supseteq K_2 \supseteq \cdots\) de compactos no vacíos de \(X\) tiene intersección no vacía.
        \item Mostrar mediante un ejemplo que la conclusión del punto (b) puede fallar si se reemplaza ``compactos'' por ``cerrados'' en un espacio no compacto.
    \end{enumerate}
\end{ejercicio}

\begin{ejercicio}{17}
    Sea \(X\) un espacio localmente compacto Hausdorff.
    \begin{enumerate}
        \item Demostrar que todo subespacio abierto de \(X\) es localmente compacto.
        \item Demostrar que todo subespacio cerrado de \(X\) es localmente compacto.
        \item Si \(K \subseteq U \subseteq X\), con \(K\) compacto y \(U\) abierto, probar que existe un abierto \(V\) tal que
        \[
            K \subseteq V \subseteq \cl{V} \subseteq U
        \]
        y \(\cl{V}\) es compacto.
    \end{enumerate}
\end{ejercicio}

\begin{ejercicio}{18}
    Sea \(C(X,Y)\) el conjunto de funciones continuas de \(X\) en \(Y\), dotado de la topología compacto-abierta.
    \begin{enumerate}
        \item Fijado \(x \in X\), demostrar que la aplicación de evaluación
        \[
            ev_x : C(X,Y) \to Y, \qquad ev_x(f) = f(x)
        \]
        es continua.
        \item Si \(K \subseteq X\) es compacto y \(F \subseteq Y\) es cerrado, demostrar que
        \[
            A_{K,F} = \{f \in C(X,Y) : f(K) \subseteq F\}
        \]
        es cerrado en \(C(X,Y)\).
        \item Analizar qué se puede afirmar sobre \(A_{K,U} = \{f \in C(X,Y) : f(K) \subseteq U\}\) cuando \(U \subseteq Y\) es abierto.
    \end{enumerate}
\end{ejercicio}

\end{document}
